\documentclass[11pt]{article}
\usepackage[margin=1in]{geometry}
\usepackage{parskip}
\usepackage{hyperref}
\usepackage{booktabs}
\usepackage{tabularx}
\usepackage{xcolor}
\usepackage[most]{tcolorbox}
\tcbuselibrary{breakable}
\tcbset{breakable}
\usepackage{enumitem}
\hypersetup{colorlinks=true, linkcolor=black, urlcolor=blue}

\newcommand{\Part}[1]{\section*{#1}}
\newcommand{\Clock}[2]{\textbf{#1} [#2]}
\newcommand{\Term}[1]{\textbf{#1}}
\newcommand{\DV}{\Term{DV}}
\newcommand{\Position}{\Term{Position}}
\newcommand{\SB}{\Term{SB}}
\newcommand{\Diamond}{\Term{Diamond}}
\newcommand{\Tier}{\Term{Tier}}

\definecolor{softgray}{RGB}{245,245,245}
\definecolor{inkblue}{RGB}{20,40,120}
\definecolor{blood}{RGB}{130,10,10}
\definecolor{forest}{RGB}{16,100,60}

\tcbset{colback=softgray, colframe=black, boxrule=0.4pt, arc=2mm}

\begin{document}

\begin{center}
{\LARGE \textbf{The Iron Crucible}}\\[4pt]
{\large A Combat-Heavy Military Campaign for \emph{Fate's Edge}}\\[2pt]
\textit{Tiers II--IV \;|\; 6--10 sessions \;|\; Siege warfare, mercenary politics, war magic}
\end{center}

\hrule
\vspace{0.5em}

\begin{quote}
\small
\textit{``The Granite Gate has never fallen. The River of Blades runs red. The Citadel of Ash waits. You are the Iron Phoenix – a company of sell‑swords, exiles, and debtors. Your contract is simple: break what cannot be broken. Your payment is a name cleared, a fortune earned, or a grave forgotten. The war machine has teeth. Let's see if yours are sharper.''}
\end{quote}

\section*{Campaign Overview}

\emph{The Iron Crucible} is a three‑part military campaign for \emph{Fate's Edge} that puts players in command of a mercenary company. Each part escalates the scale of conflict – from fortress siege to river war to city‑wide annihilation – while keeping the core mechanics of \emph{Fate's Edge} (Position, Effect, clocks, Story Beats) front and centre.

\subsection*{The Iron Phoenix}

The players are the officers of the \textbf{Iron Phoenix}, a mid‑sized mercenary company hired by a desperate coalition. The company has a reputation track, resource clocks, and a roster of lieutenants. Between missions, players manage logistics, negotiate contracts, and decide which tactics to employ.

\subsection*{The Three Crucibles}
\begin{itemize}
\item \textbf{Part I: The Granite Gate} (Tier II) – Breach an impregnable fortress using siegecraft, subterfuge, or storm.
\item \textbf{Part II: The River of Blades} (Tier III) – Control a strategic waterway, fight naval actions, and secure a bridgehead.
\item \textbf{Part III: The Citadel of Ash} (Tier IV) – Conquer a cursed city, capture or destroy magical WMDs, and face the enemy’s last stand.
\end{itemize}

\section*{Mercenary Company Rules}

The company has four clocks that represent its ability to fight. Players must manage these between sessions.

\begin{itemize}
\item \textbf{Manpower [8]} – Number of able troops. Drops from casualties. Can be replenished with coin or recruitment scenes.
\item \textbf{Supplies [6]} – Food, ammunition, and repair materials. Resupply costs silver or favours.
\item \textbf{Morale [8]} – Discipline and spirit. Low morale causes Desperate starts in battles.
\item \textbf{Reputation [6]} – Fame and trust among employers. High reputation opens better contracts and allies.
\end{itemize}

At the end of each mission (or when the party takes downtime), the GM may ask for a \textbf{Company Roll}: the player with the highest Command rolls Presence + Command (DV based on current situation). On a success, one clock of their choice advances by 1; on a failure, a clock of the GM’s choice ticks backward.

\subsection*{Company Assets}
The party can acquire assets (siege weapons, scouts, war mages, etc.) by spending \Diamond{} rewards or completing side objectives. Each asset grants a one‑time tactical bonus (e.g., +1 Position on a certain roll) or unlocks a special action.

\subsection*{Combat Position in Mass Battles}
Instead of individual positioning, squads use the same Position ladder:
\begin{itemize}
\item \textbf{Dominant} – Flanking, high ground, prepared ambush.
\item \textbf{Controlled} – Standard engagement, cover available.
\item \textbf{Desperate} – Pinned, surrounded, outnumbered.
\end{itemize}
The GM sets the Position based on the party’s tactical choices before each engagement.

\section*{Part I: The Granite Gate (Tier II)}

\subsection*{The Situation}
The Granite Gate is a fortress carved into a mountain pass, held by General Vorlag’s veteran garrison. The Iron Phoenix has been hired to open the gate for a coalition army. Three approaches are possible, each represented by a clock. The party can attempt one or more in parallel, but splitting forces risks casualties.

\begin{itemize}
\item \textbf{Main Gate Assault [8]} – Direct battering ram and ladder assault. High risk, high reward.
\item \textbf{Cliffside Scalers [6]} – Stealth climb up the sheer north face. Requires Athletics and Stealth tests.
\item \textbf{Underminers [6]} – Tunnel through the mountain’s weak rock. Requires Craft and endurance.
\end{itemize}

Once any clock fills, the party gains a foothold. To complete the siege, they must also fill the \textbf{Gatehouse Control [10]} clock. Each successful assault action ticks this clock; enemy counter‑actions can tick it backward.

\subsection*{Key Scenes}

\paragraph{The Killing Ground}
Open courtyard overlooked by murder holes. \textbf{Position: Desperate} for crossing unless the party creates cover (e.g., using siege shields, war magic, or smoke). \textbf{DV 4} for a squad to cross without casualties. On a 1, an oil pot ignites – Harm 1 to a squad.

\paragraph{Arrow Storm Passage}
A long corridor with arrow slits. \textbf{Position: Controlled} if advancing with tower shields; \textbf{Desperate} if rushing. \textbf{Clock [6]} to clear the archers. On a miss, the party loses 1 Manpower.

\paragraph{General Vorlag’s Last Stand}
The old general retreats to the inner keep. He offers a deal: safe passage for his garrison in exchange for the gate’s keys. The party must choose – honour the deal (gains Reputation, loses loot) or fight (tough battle, gains plunder but loses Morale).

\subsection*{War Magic Options (Casters/Runekeepers)}
\begin{itemize}
\item \textbf{Earthworks (Standard)} – Create a wall of packed earth (DV 4, Ritual). Grants Controlled Position for one scene.
\item \textbf{Siege Ward (Standard)} – Protect a squad from ranged attacks (gain +1 armor). Lasts one engagement.
\item \textbf{Logistics Blessing (Low)} – Create food and water for a day. Restores 1 Supply.
\end{itemize}

\subsection*{Story Beat Table (Granite Gate)}
\begin{itemize}
\item \textbf{1 SB} – A murder hole opens above a PC; test Body + Athletics or take Harm 1.
\item \textbf{2 SB} – Vorlag’s sappers collapse a tunnel. Choose: lose 2 Manpower or abort the underminer approach.
\item \textbf{3 SB} – A traitor among the coalition reveals the party’s plan. One approach clock ticks backward by 2.
\item \textbf{4+ SB} – The gate’s ancient wards activate. All magic becomes Desperate for one scene.
\end{itemize}

\section*{Part II: The River of Blades (Tier III)}

\subsection*{The Situation}
The coalition must cross the River of Blades, a fast‑flowing waterway patrolled by Admiral Kessel’s river fleet. The party must secure a bridgehead and then protect supply convoys. Three clocks track progress:

\begin{itemize}
\item \textbf{Enemy Fleet Strength [8]} – Destroy or capture enemy vessels.
\item \textbf{Bridgehead Establishment [6]} – Secure a landing zone on the far bank.
\item \textbf{Supply Line Security [6]} – Keep resupply barges safe from raiders.
\end{itemize}

The final battle is a naval engagement represented by the \textbf{Decisive River Battle [10]} clock. Each successful boarding action, ramming, or artillery strike ticks the clock; enemy counter‑attacks can reverse progress.

\subsection*{Naval Combat Actions}
Instead of individual movement, the party chooses a tactic each round:
\begin{itemize}
\item \textbf{Ramming} – Risk damage for instant boarding advantage. Test Body + Pilot (DV 4). On success, enemy vessel is locked; on failure, your ship takes Harm 1.
\item \textbf{Boarding Action} – Close combat between vessels. Use squad rules. Position depends on relative speed and size.
\item \textbf{Fire Ship} – Deploy a drifting fire barge to create an area denial zone. Any ship in the zone takes Harm 1 per round.
\item \textbf{Archery Exchange} – Ranged volley. Test Wits + Ranged (DV 3). On success, tick Enemy Fleet Strength by 1 and reduce enemy’s return fire.
\end{itemize}

\subsection*{Vessel Types (choose or earn)}

\begin{tabular}{@{}lll@{}}
\toprule
\textbf{Vessel} & \textbf{Advantages} & \textbf{Weaknesses}\\
\midrule
War Galley & +2 dice to boarding, high armour & Slow, vulnerable to fire\\
River Sloop & +1 dice to manoeuvres, fast & Low armour, small crew\\
Fire Barge & Area denial, cheap & One‑use, drifts randomly\\
Troop Transport & Carries 3 squads & Very slow, easy target\\
\bottomrule
\end{tabular}

\subsection*{Admiral Kessel’s Ambush}
The enemy admiral is a master of deception. At any point, the GM may spend 3 SB to trigger a \textbf{River Ambush}: the party’s vessel is suddenly flanked by hidden enemy boats. Position becomes Desperate for one exchange unless the party had previously scouted (Wits + Survival, DV 5).

\subsection*{Story Beat Table (River of Blades)}
\begin{itemize}
\item \textbf{1 SB} – A sudden current shift pushes the party’s boat off course. Delay 1 segment.
\item \textbf{2 SB} – A fire ship drifts into the supply convoy. Lose 1 Supply or take Harm 1 to a squad.
\item \textbf{3 SB} – Admiral Kessel personally leads a boarding party. He is a Tier III adversary with elite marines.
\item \textbf{4+ SB} – The river’s spirit awakens (angry due to the battle). It demands a sacrifice – a memory, a weapon, or a life – before letting any vessel pass.
\end{itemize}

\section*{Part III: The Citadel of Ash (Tier IV)}

\subsection*{The Situation}
The enemy has retreated to a cursed city, once a capital of a fallen empire. Three magical WMDs are hidden in the city: the \textbf{Sun‑Thrower} (laser siege weapon), the \textbf{Quake Engines} (collapse buildings), and the \textbf{Plague Cauldrons} (biological warfare). The party must capture or destroy them while controlling seven districts.

\subsection*{District Control Clock [12]}
Each district can be held by either side. The party gains 1 tick per district controlled at the end of each session. When the clock fills, the coalition declares victory. However, the enemy can launch counter‑attacks (GM spends SB) to retake districts.

\subsection*{Districts (with control values)}

\begin{tabular}{@{}lll@{}}
\toprule
\textbf{District} & \textbf{Control Value} & \textbf{Special Rule}\\
\midrule
Wizard’s Quarter & 3 & Magical hazards, teleportation circles\\
Market Square & 2 & Ambush opportunities, civilian crowds\\
Noble Heights & 2 & Defensible positions, political hostages\\
Docks & 1 & Naval support access\\
Temple District & 1 & Healing sanctuary (restore 1 Harm per scene)\\
Slums & 1 & Guerrilla warfare, hidden passages\\
Citadel & 4 (victory condition) & Final objective, heavy fortifications\\
\bottomrule
\end{tabular}

\subsection*{Magical WMD Rules}
Each weapon is a \Tier\ IV obstacle that requires a specialised skill challenge to seize.

\paragraph{Sun‑Thrower}
A massive lens that projects a beam of pure light. \textbf{Capture Clock [8]}. Each attempt requires a Wits + Craft (DV 5) to redirect the mirrors. Failure causes a burst that deals Harm 2 to the party.

\paragraph{Quake Engines}
Machines that cause building collapses. \textbf{Capture Clock [6]}. Requires Body + Athletics (DV 4) to stabilise the resonators. On a miss, the building the party is standing on collapses (Harm 1, all actions Desperate for one round).

\paragraph{Plague Cauldrons}
Vats of alchemical pestilence. \textbf{Capture Clock [6]}. Requires Spirit + Lore (DV 4) to neutralise the toxins. On a miss, the cauldron leaks – the party gains the \emph{Plague} condition (–1 die to all physical actions until they receive treatment).

\subsection*{The Bloody Baron’s Last Stand}
The enemy commander, the Bloody Baron, is a fanatic. He will not surrender. Instead, he offers a deal: he will destroy the city and all WMDs unless the party lets him escape with his elite guard. The party must choose: let him go (lose Reputation, avoid final battle) or fight him (epic confrontation, risk massive casualties).

\subsection*{Story Beat Table (Citadel of Ash)}
\begin{itemize}
\item \textbf{1 SB} – Rubble blocks a key route. Lose 1 segment of progress on a district control.
\item \textbf{2 SB} – Civilian refugees clog a district. Using area weapons would kill innocents – gain +1 Morale if you avoid them, –1 Reputation if you don’t.
\item \textbf{3 SB} – A WMD malfunctions and begins a countdown. The party has 3 rounds to disarm or evacuate the district.
\item \textbf{4+ SB} – The Hollow’s attention turns to the city. A Debt‑Walker appears and demands that the party stop using the WMDs – or it will claim the souls of everyone in the district.
\end{itemize}

\section*{Company Management Between Missions}

After each part, the party has a downtime phase. They may:

\begin{itemize}
\item \textbf{Recruit (Command + Presence, DV 3)} – Restore 1d4 Manpower.
\item \textbf{Resupply (Wits + Commerce, DV 4)} – Restore 1d4 Supplies, costing 1 Boon or a Diamond.
\item \textbf{Train (Spirit + Command, DV 4)} – Advance one squad’s specialisation (see below) or improve Morale by 1.
\item \textbf{Network (Presence + Sway, DV 3)} – Gain a new contact or a rumour about the next mission.
\end{itemize}

\subsection*{Squad Specialisations}
Once the company has Morale 6+ and has completed two missions, squads can be specialised. Choose one:

\begin{tabular}{@{}ll@{}}
\toprule
\textbf{Specialisation} & \textbf{Benefit}\\
\midrule
Shock Troops & +1 Effect on charge attacks\\
Skirmishers & Ignore difficult terrain penalties\\
Archers & Ranged attacks at +1 die\\
Engineers & +1 die to craft and siege actions\\
Medics & Reduce squad casualties by 1 level (e.g., Harm 2 → Harm 1)\\
\bottomrule
\end{tabular}

\section*{XP Progression}
XP is awarded at the end of each session as per standard \emph{Fate's Edge} guidelines, plus:

\begin{itemize}
\item Completing a part: +10 XP.
\item Capturing a WMD: +5 XP each.
\item Achieving a major moral victory (e.g., sparing civilians, refusing a dishonourable order): +3 XP.
\item Losing a squad or company asset: –2 XP (narrative loss).
\end{itemize}

\section*{NPC Adversaries (Quick Stats)}

\begin{tabular}{@{}llll@{}}
\toprule
\textbf{NPC} & \textbf{Role} & \textbf{Tier} & \textbf{Special}\\
\midrule
General Vorlag & Defensive genius & III & Fortifications gain +1 DV\\
Admiral Kessel & Naval ambusher & III & Ambush at Sea: starts one Position higher\\
Archmage Theron & War magic expert & IV & Enemy spells at +1 die\\
The Bloody Baron & Fanatic leader & IV & Morale checks at +1 DV\\
Iron Legion Captain & Elite heavy infantry & III & Armour 2, can use Shield Wall\\
\bottomrule
\end{tabular}

\section*{Running the Campaign}

\subsection*{Pacing}
\begin{itemize}
\item Part I: 2–3 sessions.
\item Part II: 2–3 sessions.
\item Part III: 3–4 sessions.
\end{itemize}
Adjust by focusing on key scenes and using SB spends to create drama.

\subsection*{Tone}
This campaign is meant to be gritty, desperate, and morally grey. Emphasise the cost of war – dead allies, ruined cities, and the weight of command. Reward creative tactics and penalise reckless bloodshed.

\subsection*{Integration with Core Fate's Edge}
All rules from the core SRD apply. Use Position/Effect for every roll. Clocks track progress and pressure. Story Beats are the GM’s tool to introduce unexpected complications. The company clocks replace detailed resource tracking.

\vfill
\begin{center}
\small{\emph{The Iron Crucible} --- A Fate's Edge Military Campaign}
\end{center}

\end{document}